\documentclass[10pt]{article}
\usepackage[margin=1in]{geometry}
\usepackage{amsmath,amssymb}
\title{Disproof of Conjecture 3486: a separating pair on seven vertices}
\author{Local verified working submission}
\date{October 3, 2026}
% Compact consistent title for a short mathematical note.
\makeatletter
\renewcommand{\maketitle}{\begin{center}
{\Large\bfseries\@title\par}\vspace{0.6em}
{\small\@author\par\@date}
\end{center}\vspace{0.5em}}
\makeatother
\setlength{\emergencystretch}{3em}
\usepackage{url}
\begin{document}
\maketitle

\paragraph{Claim being disproved.}
The conjecture asserts that the smallest pair of nonisomorphic
four-critical graphs with identical single-vertex-deletion chromatic
spectra has thirteen vertices. We exhibit such a pair on seven vertices.
Here \emph{four-critical} has its usual strong meaning: the graph has
chromatic number four and every proper subgraph is three-colorable.

\paragraph{The graphs.}
Both graphs have vertex set $\{0,\ldots,6\}$. Let
\begin{align*}
E(G)={}&\{02,03,05,06,12,13,14,23,45,46,56\},\\
E(H)={}&\{01,02,04,05,12,13,15,23,24,36,46,56\}.
\end{align*}
The first graph is a Haj\'os join of two copies of $K_4$:
delete edges $01$ and $04$ in the copies on $\{0,1,2,3\}$
and $\{0,4,5,6\}$, and insert edge $14$.
The second is the Mycielski graph of $K_3$: vertices $0,1,2$ form
the original triangle, vertices $3,4,5$ copy their respective
neighborhoods, and vertex $6$ is adjacent to $3,4,5$.

\paragraph{Chromatic number four.}
In any three-coloring of the subgraph of $G$ on $\{0,1,2,3\}$,
vertices $2,3$ have different colors and both $0,1$ are adjacent
to both of them. Thus $0,1$ must have the same third color.
The same argument on $\{0,4,5,6\}$ forces $0,4$ to have the same
color. This contradicts edge $14$.
For $H$, the original triangle uses all three colors.
Vertex $3$ must then have the color of $0$, vertex $4$ the color
of $1$, and vertex $5$ the color of $2$. Vertex $6$ sees all three
colors, again a contradiction.
In contrast, the following vectors, in vertex order $0,\ldots,6$,
are proper four-colorings:
\[
 c_G=(0,0,1,2,1,2,3),\qquad
 c_H=(0,1,2,0,1,2,3).
\]
Hence $\chi(G)=\chi(H)=4$.

\paragraph{Certificates for every single deletion.}
Each vector below uses only colors $0,1,2$. In the vertex-deletion
table the entry at the deleted vertex is ignored (denoted by $*$).
The listed edge sets make every certificate directly checkable.
\begin{center}
\begin{tabular}{c|c|c}
Deleted vertex & coloring of $G-x$ & coloring of $H-x$\\ \hline
$0$ & $(*,0,1,2,1,0,2)$ & $(*,0,1,2,0,2,1)$ \\
$1$ & $(0,*,1,2,0,1,2)$ & $(0,*,1,0,2,2,1)$ \\
$2$ & $(0,1,*,2,0,1,2)$ & $(0,1,*,0,2,2,1)$ \\
$3$ & $(0,1,2,*,0,1,2)$ & $(0,1,2,*,1,2,0)$ \\
$4$ & $(0,0,1,2,*,1,2)$ & $(0,1,2,0,*,2,1)$ \\
$5$ & $(0,0,1,2,1,*,2)$ & $(0,1,2,0,1,*,2)$ \\
$6$ & $(0,0,1,2,1,2,*)$ & $(0,1,2,0,1,2,*)$ \\
\end{tabular}
\end{center}

\begin{center}
\begin{tabular}{c|c}
Deleted edge of $G$ & coloring of $G-e$\\ \hline
$02$ & $(0,1,0,2,0,1,2)$ \\
$03$ & $(0,1,2,0,0,1,2)$ \\
$05$ & $(0,0,1,2,1,0,2)$ \\
$06$ & $(0,0,1,2,1,2,0)$ \\
$12$ & $(0,1,1,2,0,1,2)$ \\
$13$ & $(0,1,2,1,0,1,2)$ \\
$14$ & $(0,0,1,2,0,1,2)$ \\
$23$ & $(0,1,2,2,0,1,2)$ \\
$45$ & $(0,0,1,2,1,1,2)$ \\
$46$ & $(0,0,1,2,1,2,1)$ \\
$56$ & $(0,0,1,2,1,2,2)$ \\
\end{tabular}\qquad
\begin{tabular}{c|c}
Deleted edge of $H$ & coloring of $H-e$\\ \hline
$01$ & $(0,0,1,2,2,1,0)$ \\
$02$ & $(0,1,0,2,1,2,0)$ \\
$04$ & $(0,1,2,0,0,2,1)$ \\
$05$ & $(0,1,2,0,1,0,2)$ \\
$12$ & $(0,1,1,0,2,2,1)$ \\
$13$ & $(0,1,2,1,1,2,0)$ \\
$15$ & $(0,1,2,0,1,1,2)$ \\
$23$ & $(0,1,2,2,1,2,0)$ \\
$24$ & $(0,1,2,0,2,2,1)$ \\
$36$ & $(0,1,2,0,1,2,0)$ \\
$46$ & $(0,1,2,0,1,2,1)$ \\
$56$ & $(0,1,2,0,1,2,2)$ \\
\end{tabular}
\end{center}

Every proper subgraph omits either a vertex or, if it has every
vertex, an edge. Restrict the corresponding coloring from these
tables to that subgraph. This proves that every proper subgraph
of either graph is three-colorable. Both graphs are therefore
four-critical in the strong sense.

\paragraph{Identical vertex-deletion spectra.}
The graph $G$ contains the two vertex-disjoint triangles $123$
and $456$, so at least one survives every single vertex deletion.
For $H$, after deleting $0,1,2$ respectively, the triangles
$123$, $024$, $015$ survive; after deleting any of $3,4,5,6$,
the triangle $012$ survives. Thus all these vertex-deleted graphs
require at least three colors. Combined with the table, this gives
\[
 (\chi(G-x))_{x=0}^6=(3,3,3,3,3,3,3)
   =(\chi(H-x))_{x=0}^6.
\]

\paragraph{Nonisomorphism and contradiction.}
There are eleven edges in $G$ and twelve in $H$, so the graphs
are nonisomorphic. An additional structural distinction, used
in Lean, is that $G$ has two vertex-disjoint triangles and $H$
does not. Indeed, vertex $6$ belongs to no triangle and
$\{3,4,5\}$ is independent. Every triangle of $H$ therefore
uses at least two of the three vertices $\{0,1,2\}$, so any
two such triangles intersect.
We have exhibited a valid pair at order $7<13$, disproving
the claimed minimum thirteen. We do not claim to determine
the true minimum order.

\paragraph{Lean verification and semantic correspondence.}
The self-contained file \texttt{Main.lean} imports only
\texttt{Std}. It defines the actual graphs by the displayed edge lists.
The definitions \texttt{Proper}, \texttt{ProperDeleted},
and \texttt{ProperEdgeDeleted} express proper colorings before
deletion and after deleting an arbitrary vertex or edge.
All color vectors in the tables are verified for every edge;
exhaustive kernel-checked tests also exclude every three-coloring
of the original graphs and every two-coloring of a vertex deletion.
The theorem \texttt{tuple\_eta} proves that the finite tuple
enumeration covers every coloring function, so these tests
exclude arbitrary colorings, not only a selected sample.

The theorem \path{proper_subgraph_bridge} formally restricts
the single-deletion colorings to every proper subgraph.
Consequently \texttt{G\_four\_critical} and
\texttt{H\_four\_critical} establish strong four-criticality.
The theorem \path{same_deletion_spectrum} proves the exact
chromatic number three, including minimality, for every vertex
deletion of both graphs.
The theorem \texttt{graphs\_not\_isomorphic} excludes an actual
adjacency-preserving vertex bijection by the disjoint-triangle
invariant.
Finally, \path{not_minimum_thirteen_consequence} refutes the
necessary order-seven consequence of the asserted minimum thirteen.

\paragraph{Reproduction and trust.}
Run \texttt{lean Main.lean} with Lean 4.19.0.
Finite checks use kernel-checked \texttt{decide}; no
\texttt{native\_decide}, unproved placeholder, or additional axiom
is used. The file prints axiom dependencies for its main
nonexistence, criticality, spectrum, and nonisomorphism results.
\end{document}

